%\documentclass[12pt,oneside]{amsart}
\documentclass{scrartcl}

%\documentclass{article}

\usepackage{amsthm,amsmath,amssymb}
\usepackage[numbers]{natbib}

%\usepackage[colorlinks]{hyperref}
\usepackage{hyperref}
\hypersetup{
    colorlinks,%
    citecolor=black,%
    filecolor=black,%
    linkcolor=black,%
    urlcolor=black
}
\usepackage{hypernat}
\usepackage[top=2.5cm, bottom=2.5cm, left=2.8cm,
right=2.8cm]{geometry} 

%\setlength{\parskip}{0pt}
%\setlength{\parsep}{0pt}
%\setlength{\headsep}{0pt}
%\setlength{\topskip}{0pt}
%\setlength{\topmargin}{0pt}
%\setlength{\topsep}{0pt}
%\setlength{\partopsep}{0pt}

\linespread{1}

%\usepackage{marvosym}

%\textwidth = 6.5 in \textheight = 9.2 in \oddsidemargin = 0.0 in
%\evensidemargin = 0.0 in \topmargin = -0.0 in \headheight = 0.0 in
%\headsep = 0.0in \parskip = 0.0in \parindent = 0.0in
%\def\baselinestretch{0.871}


\newtheorem{theorem}{Theorem}[section]
\newtheorem{proposition}[theorem]{Proposition}
\newtheorem{problem}[theorem]{Problem}
\newtheorem{fact}[theorem]{Fact}
\newtheorem{lemma}[theorem]{Lemma}
\newtheorem{defn}[theorem]{Definition}
\newtheorem{corollary}[theorem]{Corollary}
\newtheorem{remark}{Remark}[section]
\newtheorem{example}[theorem]{Example}

\newcommand{\E}{{\mathbb E}}
\newcommand{\se}{{\mathcal{E}}}
\newcommand{\R}{{\mathbb R}}
\renewcommand{\P}{{\mathbb P}}
\newcommand{\ac}{{\mathcal{A}}}
\newcommand{\A}{{\mathcal{A}}}
\newcommand{\B}{{\mathcal{B}}}
\newcommand{\C}{{\mathcal{C}}}
\newcommand{\M}{{\mathcal{M}}}
\newcommand{\Mf}{{\mathfrak{M}}}
\newcommand{\T}{{\mathcal{T}}}
\newcommand{\Z}{{\mathcal{Z}}}
\newcommand{\K}{{\mathcal{K}}}
\newcommand{\lc}{{\mathcal{L}}}
\newcommand{\Ps}{{\mathcal{P}}}
\newcommand{\G}{{\mathcal{G}}}
\newcommand{\pa}{{\mathring{p}}}
\newcommand{\F}{{\cal F}}
\newcommand{\samp}{{\mathcal{X}}}
\newcommand{\X}{{\mathcal{X}}}
\newcommand{\hi}{{\hat{\Phi}}}
\newcommand{\data}{{\text{data}}}

\newcounter{rcnt}[section]
\renewcommand{\thercnt}{(\roman{rcnt})}

\def\qt#1{\qquad\text{#1}}

\def\argmin{\mathop{\rm argmin}}
\def\argmax{\mathop{\rm argmax}}


\setlength{\parskip}{1.4 \medskipamount}
%\sloppy
%\linespread{1.3}

\begin{document}

\title{STAT 153 \& 248 - Time Series  \\
  Lecture Three} 
\subtitle{Fall 2026, UC Berkeley}
\author{Aditya Guntuboyina}

\date{03 September 2026}

\maketitle

We shall look at the math behind the linear regression
calculations. We will mostly focus on the Bayesian approach which is
conceptually simpler and generalizes easily to more complicated
models. We shall look at the main frequentist ideas in the next
lecture.

In order to motivate the main ideas behind Bayesian inference for
linear regression, we shall start with some simple inference problems
first before looking at linear regression.

\section{Simple Question 1}

Suppose $a$ and $b$ are two numbers and all we know is that $a + b =
7$. Based on this information, what can we say about $a$?

There is clearly not enough information to determine $a$. For every choice of
$a$, we can choose $b=7-a$. Thus there are infinitely many possible
pairs $(a,b)$.

\section{Simple Question 2}

Suppose $a$ and $b$ are two numbers such that $a + b = 7$. Suppose $b$
is small in magnitude. What then can be said about $a$?

Here is a way of arriving at a reasonable solution to this problem
using probability theory. Assume that $a$ and $b$ are independent
with
\begin{align}\label{md2}
  a \sim \operatorname{Unif}(-C,C) ~~ \text{ and } ~~
  b \sim N(0,\sigma^2)
\end{align}
where $C$ is very large, and $\sigma$ is a fixed quantity that
characterizes smallness of $b$. Depending on what we understand by
smallness of $b$, we can set a value of $\sigma$ (e.g., $\sigma =
0.5$).

Using the \textit{modeling} assumption \eqref{md2} and the available
\textit{data} that $a + b = 7$, our uncertainty in $a$ will be
captured by the probability distribution:
\begin{align*}
  a \mid a + b = 7. 
\end{align*}
Here is a simple way of calculating this conditional distribution. Let
$y = a + b$ so that we need to calculate the conditional density of $a
\mid y = 7$. The conditional distribution of $y \mid a$ is given by
$N(a, \sigma^2)$ because $y = a + b$,  and $b$ is distributed as $N(0,
\sigma^2)$  independently of $a$. Therefore by Bayes rule:
\begin{align*}
  f_{a \mid y = 7}(a) &\propto f_{y \mid a}(7) f_{a}(a) \\
  &= \frac{1}{\sigma \sqrt{2 \pi}} \exp \left(-\frac{(7 - a)^2}{2
    \sigma^2} \right) \times \frac{I\{-C < a < C\}}{2C} \\
  &\propto \exp \left(-\frac{(7 - a)^2}{2
    \sigma^2} \right) I\{-C < a < C\}. 
\end{align*}
This is the normal distribution $N(7, \sigma^2)$ truncated to the
interval $(-C, C)$. Because we shall take $C$ to be very large, the
restriction to $(-C, C)$ has almost no effect and therefore,
\begin{align*}
  a\mid a+b=7 &\approx N(7,\sigma^2). 
\end{align*}
Thus under our probabilistic assumptions and the given information,
our uncertainty about $a$ is captured by the normal distribution
centered at 7 and variance $\sigma^2$. The amount of uncertainty is
controlled by $\sigma^2$. 

\section{Simple Question 3}

Now suppose we have three numbers $a, b_1, b_2$. Assume that both
$b_1$ and $b_2$ are small in magnitude. Also suppose that 
\begin{align*}
  a+b_1 = 7 ~~ \text{ and } ~~ a+b_2 &= 10. 
\end{align*}
What can then be said about $a$?

Note now that there are three unknown numbers and two pieces of
information. So this problem is also logically unanswerable without
further assumptions. We shall make probabilistic assumptions on the
unknown $a, b_1, b_2$. Specifically, we assume that 
\begin{align}\label{md3}
  a \sim \operatorname{Unif}(-C,C),  ~~~ \text{ and } ~~~  b_1,b_2\mid\sigma
  \overset{\mathrm{iid}}{\sim}N(0,\sigma^2) ~~~ \text{ and } ~~~
  \log\sigma \sim \operatorname{Unif}(-C,C)
\end{align}
where, as before, $C$ is a very large constant. The idea behind this
modeling assumption is:
\begin{enumerate}
\item We don't want to assume anything specific about $a$
\item Because $b_1, b_2$ are given to be small in magnitude, we assume
  a normal distribution centered at 0 for them with their magnitude
  controlled by a parameter $\sigma$. This also assumes they have the
  same scale.
\item The common scale $\sigma$ of $b_1, b_2$ is also unknown. We
  don't want to assume anything specific about this scale so we use
  the uniform $(-C, C)$ prior on $\log \sigma$. Because $\sigma$ is
  always positive, $\log \sigma$ makes sense and it takes values that
  can be both positive and negative, and this makes the
  $\text{Unif}(-C, C)$ prior reasonable for $\log \sigma$. 
\end{enumerate}
The assumption $\log \sigma \sim \text{unif}(-C, C)$ leads to the
following density for $\sigma$ (by the change of variables formula): 
The last assumption says that we do not know the order of magnitude of
$\sigma$. By a change of variables, it gives the density
\begin{align*}
  f_\sigma(\sigma)
  = f_{\log\sigma}(\log\sigma)
  \left|\frac{d}{d\sigma}\log\sigma\right| = f_{\log \sigma}(\log
  \sigma) \frac{1}{\sigma} = \frac{I\{-C < \log \sigma < C\}}{2 C
  \sigma}. 
\end{align*}
With these assumptions, we need to calculate the conditional
distribution of:
\begin{align*}
  a \mid a+ b_1 = 7, a + b_2 = 10. 
\end{align*}
This determines our uncertainty about $a$ given the information $a +
b_1 = 7$ and $a + b_2 = 10$. To compute the above conditional, let us
use the notation:
\begin{align*}
  y_1 = a + b_1 ~~~ \text{ and } ~~~ y_2 = a + b_2, 
\end{align*}
so that we need to calculate the density of:
\begin{align}\label{req_cond}
  a \mid y_1 = 7, y_2 = 10. 
\end{align}
One complication in calculating this density is the variable $\sigma$
which is present in the modeling assumption but is not present in this
conditional density. A simple way of dealing with this is to first
compute the joint conditional density
\begin{align}\label{req_cond_sig}
  a, \sigma \mid y_1 = 7, y_2 = 10
\end{align}
and then to integrate over $\sigma$ to obtain the required conditional
density \eqref{req_cond}. To calculate \eqref{req_cond_sig}, we use
the Bayes rule in the following way: (note that $y_1, y_2 \mid a,
\sigma \sim N(a, \sigma^2)$)
\begin{align*}
  f_{a, \sigma \mid y_1, y_2}(a, \sigma) &\propto
                                                         f_{y_1, y_2
                                                         \mid a,
                                                         \sigma}(y_1,  y_2)
                                                         f_{a,
                                                         \sigma}(a,
                                                         \sigma) \\
  &= \left\{\prod_{j=1}^2
  \frac{1}{\sqrt{2\pi}\sigma}
  \exp\left(-\frac{(y_j-a)^2}{2\sigma^2}\right)\right\} f_a(a)
    f_{\sigma}(\sigma) \\
  &= \left\{\prod_{j=1}^2
  \frac{1}{\sqrt{2\pi}\sigma}
  \exp\left(-\frac{(y_j-a)^2}{2\sigma^2}\right)\right\} \frac{I\{-C <
    a < C\}}{2 C} \frac{I\{-C < \log \sigma < C\}}{2 C \sigma} \\
  &\propto \sigma^{-3} \exp \left(-\frac{(y_1 - a)^2 + (y_2 - a)^2}{2
    \sigma^2} \right) I\{-C < a < C\} I\{e^{-C} < \sigma < e^C\}
\end{align*}
where, in the last step above, we ignored the terms involving
$\sqrt{2\pi}$ and $C$ in the proportional constant. As $C$ is
very large, we shall also ignore the first indicator function above because
it will be 1 for all reasonable values of $a$, and replace the second
indicator function by $I\{\sigma > 0\}$ because it will be 1 for all
reasonable values of $\sigma > 0$. This gives
\begin{align*}
  f_{a, \sigma \mid y_1, y_2}(a, \sigma) \propto \sigma^{-3} \exp \left(-\frac{(y_1 - a)^2 + (y_2 - a)^2}{2
    \sigma^2} \right) I\{\sigma > 0\}. 
\end{align*}
This conditional density reflects our knowledge of both $a$ and
$\sigma$ given the information on $y_1$ and $y_2$. Since the question
only cares about $a$ (and not $\sigma$), we can obtain
\eqref{req_cond} by integrating the above with respect to
$\sigma$. This gives:
\begin{align*}
  f_{a \mid y_1, y_2}(a) &= \int f_{a, \sigma \mid y_1, y_2}(a, \sigma)
                           d\sigma \propto \int_0^{\infty} \sigma^{-3} \exp \left(-\frac{(y_1 - a)^2 + (y_2 - a)^2}{2
                           \sigma^2} \right) d \sigma.
\end{align*}
Let
\begin{align*}
  S(a) &= (y_1-a)^2+(y_2-a)^2, 
\end{align*}
and we can call $S(a)$ the sum of squares function. Then
\begin{align*}
  f_{a \mid y_1, y_2}(a) \propto \int_0^{\infty} \sigma^{-3} \exp \left(-\frac{S(a)}{2
                           \sigma^2} \right) d \sigma.
\end{align*}
We use the change of variable:
\begin{align}\label{cov}
  \sigma = t \sqrt{S(a)} ~~~ \text{ so that } ~~~ d \sigma =
  \sqrt{S(a)}~dt. 
\end{align}
Then
\begin{align*}
  f_{a \mid y_1, y_2}(a) &\propto \int_0^{\infty} \sigma^{-3} \exp \left(-\frac{S(a)}{2
                           \sigma^2} \right) d \sigma \\
  &= \int_0^{\infty} t^{-3} S(a)^{-3/2} \exp \left(-\frac{1}{2 t^2}
    \right) \sqrt{S(a)} dt \\
  &= \frac{1}{S(a)} \int_0^{\infty} t^{-3} \exp \left(-\frac{1}{2t^2}
    \right) dt \\
  &\propto \frac{1}{S(a)}, 
\end{align*}
where, in the last step, we ignored the integral term because it is a
constant (i.e., does not depend on $a$). We therefore proved
\begin{align*}
  f_{a\mid y_1,y_2}(a) \propto \frac{1}{(y_1-a)^2+(y_2-a)^2}.
\end{align*}
It turns out that this is a Cauchy density. To see this, first
complete the square and rewrite $S(a)$ as: 
\begin{align*}
  (y_1-a)^2+(y_2-a)^2
  =2\left(a-\frac{y_1+y_2}{2}\right)^2
    +\frac{(y_1-y_2)^2}{2} =2\left\{
    \left(a-\frac{y_1+y_2}{2}\right)^2
    +\left(\frac{|y_1-y_2|}{2}\right)^2
    \right\}
\end{align*}
so that
\begin{align}\label{tcau_1}
  f_{a \mid y_1, y_2}(a) \propto \frac{1}{    \left(a-\frac{y_1+y_2}{2}\right)^2
    +\left(\frac{|y_1-y_2|}{2}\right)^2}. 
\end{align}
The density of a Cauchy random variable with location $\mu$ and scale
$\tau$ is
\begin{align} \label{tcau_2}
 \text{Density of Cauchy}(\mu, \tau) = \frac{1}{\pi\tau} \frac{1}{1+\left((x-\mu)/\tau\right)^2}
  \propto \frac{1}{\tau^2+(x-\mu)^2}. 
\end{align}
Comparing \eqref{tcau_1} and \eqref{tcau_2}, we get 
\begin{align*}
  a\mid y_1,y_2
  \sim \operatorname{Cauchy}\left(
  \frac{y_1+y_2}{2},\frac{|y_1-y_2|}{2}\right).
\end{align*}
For $y_1=7$ and $y_2=10$, this becomes
\begin{align*}
  a\mid y_1=7,y_2=10
  &\sim \operatorname{Cauchy}(8.5,1.5).
\end{align*}
It is quite remarkable that we obtain a Cauchy density to capture the
uncertainty about $a$ given the observations. This is basically a
consequence of the modeling assumptions \eqref{md3}. 

\section{Question 4: Several Measurements of One Quantity}

Suppose a scientist makes $n=6$ measurements of an unknown physical
quantity $\theta$ and obtains 
\begin{align*}
  y_1=26.6, ~~ y_2=38.5, ~~ y_3 = 34.4, ~~ y_4 = 34.0, ~~ y_5 = 31.0,
  ~~
  y_6
  =
  23.6. 
\end{align*}
Based on these observations, what can be inferred about $\theta$?

We consider each observation $y_i$ as comprised of
$\theta$ and an 'error' or 'noise' term $\epsilon_i$:
\begin{align*}
  y_i = \theta + \epsilon_i \qt{for $i = 1, \dots, n$}. 
\end{align*}
This means that we are trying to write the $n$ observations $y_1,
\dots, y_n$ in terms of $n+1$ unknowns $\theta$ and $\epsilon_1,
\dots, \epsilon_n$. On these $n+1$ unknowns, we make the following
modeling assumptions: 
\begin{align*}
  \theta \sim \operatorname{Unif}(-C,C) ~~ \text{ and } ~~ 
  \epsilon_1,\ldots,\epsilon_n\mid\sigma
  \overset{\mathrm{iid}}{\sim}N(0,\sigma^2) ~~ \text{ and } ~~ 
  \log\sigma \sim \operatorname{Unif}(-C,C).
\end{align*}
This assumption is the same as that used in the previous section, the
only difference is that $n$ is now 6 as opposed to 2, and also we are
using different notation ($\theta$ in place of $a$, and $\epsilon_1,
\dots, \epsilon_n$ in place of $b_1, b_2$).

Our goal is to find the conditional density of $\theta$ given $y_1,
\dots, y_n$. As in question 3, we first find the joint conditional
density of $(\theta, \sigma)$ given $y_1, \dots, y_n$, and then we
integrate out $\sigma$ to calculate $\theta \mid y_1, \dots, y_n$. By
Bayes rule,
\begin{align*}
  f_{\theta, \sigma \mid y_1, \dots, y_n}(\theta, \sigma) &\propto
                                                         f_{y_1,
                                                            \dots, y_n
                                                         \mid \theta,
                                                         \sigma}(y_1,
                                                            \dots,  y_n)
                                                         f_{\theta,
                                                         \sigma}(\theta,
                                                         \sigma) \\
  &= \left\{\prod_{j=1}^n
  \frac{1}{\sqrt{2\pi}\sigma}
  \exp\left(-\frac{(y_j-\theta)^2}{2\sigma^2}\right)\right\} f_{\theta}(\theta)
    f_{\sigma}(\sigma) \\
  &= \left\{\prod_{j=1}^n
  \frac{1}{\sqrt{2\pi}\sigma}
  \exp\left(-\frac{(y_j-\theta)^2}{2\sigma^2}\right)\right\} \frac{I\{-C <
    \theta < C\}}{2 C} \frac{I\{-C < \log \sigma < C\}}{2 C \sigma} \\
  &\propto \sigma^{-n-1} \exp \left(-\frac{(y_1 - \theta)^2 + \dots +
    (y_n - \theta)^2}{2 
    \sigma^2} \right) I\{-C < \theta < C\} I\{e^{-C} < \sigma <
    e^C\}. 
\end{align*}
Ignoring the indicators, we get 
\begin{align*}
  f_{\theta, \sigma \mid y_1, \dots, y_n}(\theta, \sigma) \propto
  \sigma^{-n-1} \exp \left(-\frac{(y_1 - \theta)^2 + \dots + (y_n -
  \theta)^2}{2  \sigma^2} \right) I\{\sigma > 0\}. 
\end{align*}
Integrating the above with respect to $\sigma$, we get 
\begin{align*}
  f_{\theta \mid y_1, \dots, y_n}(\theta) &= \int f_{\theta, \sigma \mid
                                            y_1, \dots, y_n}(\theta, \sigma) 
                           d\sigma \propto \int_0^{\infty}
                                            \sigma^{-n-1} \exp
                                            \left(-\frac{(y_1 - \theta)^2 +
                                            \dots +
                                            (y_n - \theta)^2}{2 
                           \sigma^2} \right) d \sigma.
\end{align*}
Let
\begin{align*}
  S(\theta) &= (y_1-\theta)^2+ \dots +(y_n-\theta)^2
\end{align*}
be the sum of squares function. Then
\begin{align*}
  f_{\theta \mid y_1, \dots, y_n}(\theta) \propto \int_0^{\infty}
  \sigma^{-n-1} \exp \left(-\frac{S(\theta)}{2 
                           \sigma^2} \right) d \sigma.
\end{align*}
Using the same change of variable as in \eqref{cov}, we deduce
\begin{align}\label{tt_1}
  f_{\theta \mid y_1, \dots, y_n}(\theta) \propto \left(
  \frac{1}{S(\theta)} \right)^{n/2}. 
\end{align}
We will see next week that this is a $t$-density with $n-1$ degrees of
freedom. Note that the value of $\theta$ that  minimizes $S(\theta)$
is the sample mean:
\begin{align*}
  \bar{y} = \frac{y_1 + \dots + y_n}{n}. 
\end{align*}
Because $1/S(\theta)$ appears in \eqref{tt_1}, it means that the
conditional density of $\theta$ given $y_1, \dots, y_n$ is maximized
at the sample mean.


\section{Linear Regression with Time as a Covariate}

Suppose $y_1, \dots, y_n$ is a time series. We want to decompose this
time series as
\begin{align*}
  y_i = \beta_0 + \beta_1 x_i + \epsilon_i
\end{align*}
where $x_i$ is known. The simplest example is $x_i = i$; this
corresponds to regression with time as covariate. The analysis also
works when $x_i = y_{i-1}$ as in lagged regression but we shall study
this later. There are $n+2$ unknowns here: $\beta_0, \beta_1$ and
$\epsilon_1, \dots, \epsilon_n$. The goal is to figure them out from
the $n$ known observations $y_1, \dots, y_n$.

We make the following modeling assumptions:
\begin{align*}
  \beta_0 \sim \text{Unif}(-C, C) ~~ \beta_1 \sim \text{Unif}(-C, C)
  ~~ \epsilon_1, \dots, \epsilon_n \mid \sigma
  \overset{\text{i.i.d}}{\sim} N(0, \sigma^2) ~~ \log \sigma \sim
  \text{Unif}(-C, C). 
\end{align*}
Our interest is usually in $\beta_0, \beta_1$. So we aim to calculate
the conditional density of $(\beta_0, \beta_1)$ given $y_1, \dots,
y_n$. We do this by first calculating the density of $(\beta_0,
\beta_1, \sigma)$ given $y_1, \dots, y_n$, and then integrating out
$\sigma$. By Bayes rule, 
\begin{align*}
  &f_{\beta_0,\beta_1,\sigma\mid y_1,\ldots,y_n}
  (\beta_0,\beta_1,\sigma) \\
  &\qquad\propto
  \left\{\prod_{i=1}^n
  \frac{1}{\sqrt{2\pi}\sigma}
  \exp\left[-\frac{(y_i-\beta_0-\beta_1x_i)^2}{2\sigma^2}\right]
  \right\}
  \frac{I\{-C<\beta_0<C\}}{2C}
  \frac{I\{-C<\beta_1<C\}}{2C} \\
  &\qquad\quad\times
  \frac{ I\{-C<\log\sigma<C\}}{2C\sigma} \\
  &\qquad\propto \sigma^{-(n+1)}
  \exp\left\{-\frac{S(\beta_0,\beta_1)}{2\sigma^2}\right\} I\{\sigma >
    0\},
\end{align*}
where 
\begin{align*}
  S(\beta_0,\beta_1) =\sum_{i=1}^n(y_i-\beta_0-\beta_1x_i)^2
\end{align*}
is the sum of squares, and we have ignored the indicator $I\{-C <
\beta_0 < C\}$, $I\{-C < 
\beta_1 < C\}$ and replaced the indicator $I\{e^{-C} < \sigma < e^C\}$
by $I\{\sigma > 0\}$. 

As in the previous two sections, integrating out $\sigma$ gives 
\begin{align*}
  f_{\beta_0,\beta_1\mid y_1,\ldots,y_n}
  (\beta_0,\beta_1\mid y_1,\ldots,y_n)
  &\propto
  \left\{\frac{1}{S(\beta_0,\beta_1)}\right\}^{n/2}.
\end{align*}
We will discuss in the next lecture that this is a multivariate
$t$-distribution with $n-2$ degrees of freedom. It is centered at the
least squares estimators of $\beta_0$ and $\beta_1$. 




\end{document}

%%% Local Variables:
%%% mode: latex
%%% TeX-master: t
%%% End:
